\section{Integral of a Rational Function}
%
In this section, we will present a method to explicitly calculate the integral
of any rational function.
\index{integral!of a rational function}%
\index{primitive!of a rational function}%

\begin{definition}[rational function]
\mymargin{rational function}%
\index{function!rational}%
\index{rational!function}%
A function $f\colon A \subset \RR \to \RR$
is said to be \emph{rational} if it can be expressed in the form
\[
  f(x) = \frac{P(x)}{Q(x)}
\]
where $P$ and $Q$ are polynomial functions with real coefficients.
\end{definition}

The method we are going to present can be used "blindly" because once the
decomposition is found, there is no need to know why the method works (and why
it always works).
The proof of why this method works is a non-trivial exercise in linear algebra
to which we will dedicate the rest of this chapter because, all in all, we
believe it is always useful to have an understanding of the tools being used.

\subsection{Examples with Denominator of First and Second Degree}

Before even stating the theorem and proving it, we want to present the method
through examples.
This is probably the simplest way to understand how to proceed.

\begin{example}[denominator of degree $1$]
  Calculate 
  \[
  \int \frac{x^3+1}{x-1}\, dx.  
  \]
\end{example}
\begin{proof}[Solution.]
We are faced with the integral of a rational function with a denominator of
degree 1 and a numerator of degree 3.
In general, when the degree of the numerator is not less than that of the
denominator, one must first perform the division between polynomials
(see~\ref{th:divisione_polinomi}). We find:
\[
\frac{x^3+1}{x-1} = x^2 + x + 1 + \frac{2}{x-1}.
\]
Calculating the integral of the quotient polynomial is always trivial:
\[
\int \enclose{x^2+x+1}  \, dx \ni \frac{x^3}{3} + \frac{x^2}{2} + x.
\]
Since the remainder in the division between polynomials has a degree lower than
the denominator, we always reduce to the integral of a rational function where
the numerator has a degree lower than the denominator.
If the denominator has degree $1$, as in this example, the numerator will always
have degree $0$ and thus is constant.
Therefore, up to multiplicative constants, we always reduce to the integral of a
function of the type $\frac{1}{x+c}$ whose primitive is a logarithm:
\[
 \int \frac{2}{x-1}\, dx \ni 2 \ln \abs{x-1} = \ln (x-1)^2
\]
In conclusion, one of the sought primitives is 
\[
  \int \frac{x^3+1}{x-1}\, dx \ni \frac{x^3}{3} + \frac{x^2}{2} + x + \ln (x-1)^2.
\]
\end{proof}

\begin{example}[denominator of degree 2 with distinct real roots]
Calculate 
\[
  \int \frac{x+1}{x^2-x}\, dx.
\]
\end{example}
\begin{proof}[Solution.]
The polynomial in the denominator factors as $x^2 -x = x (x-1)$. 
Then the idea is to think that the second-degree denominator was obtained as a
common denominator in a sum of rational functions with first-degree
denominators:
\begin{equation}\label{eq:541177}
  \frac{x+1}{x(x-1)} \stackrel?= \frac{A}{x} + \frac{B}{x-1}  
\end{equation}
with $A,B\in \RR$ constants to be determined. 
This is called \emph{partial fraction decomposition}.
\index{decomposition!into partial fractions}%
The constants can be determined by working backwards:
\[
  \frac{A}{x} + \frac{B}{x-1}    
  = \frac{A(x-1) + Bx}{x(x-1)} \stackrel!= \frac{x+1}{x(x-1)}.
\]
For the equality to hold, we must have 
\[
  A(x-1) + Bx = (A+B)x - A \stackrel!= x+1  
\]
and thus $A+B = 1$ and $-A=1$ from which $A=-1$ and $B=2$.
%
\mynote{\textbf{method of residues.} 
\index{residues!method of}%
\index{method!of residues}%
The coefficients $A$ and $B$ can also be found
in the following way. 
If I multiply equation~\eqref{eq:541177}
by $x$ and then let $x$ tend to $0$,
the term $B$ on the right-hand side cancels 
and I immediately obtain $A=-1$. 
If instead I multiply the equation by $x-1$ and then let $x$ tend to $1$,
the term $A$ on the right-hand side cancels
and I immediately obtain $B=2$.}
%
Thus 
\[
  \frac{x+1}{x(x-1)} = \frac{-1}{x} + \frac{2}{x-1}  
\]
and the primitive is a sum of logarithms:
\[
  \int \frac{x+1}{x^2-x}\, dx
  \ni -\ln\abs{x} + 2 \ln \abs{x-1} 
  = \ln \frac{(x-1)^2}{\abs{x}}.
\]
\end{proof}

\begin{example}[denominator of second degree with a double root]
Calculate 
\[
  \int \frac{x-1}{x^2+2x+1}\, dx
\]
\end{example}
\begin{proof}[Solution.]
The polynomial in the denominator is a perfect square $x^2+2x+1=(x+1)^2$.
In this case, the partial fraction decomposition cannot have 
the two denominators equal but one must raise a denominator to 
the square:
\begin{equation}\label{eq:2888323}
\frac{x-1}{(x+1)^2} \stackrel?= \frac{A}{x+1} + \frac{B}{(x+1)^2}.  
\end{equation}
For the equality to hold, we must have
\[
\frac{A}{x+1} + \frac{B}{(x+1)^2} = \frac{A(x+1) + B}{(x+1)^2}  
\stackrel!=\frac{x-1}{(x+1)^2}. 
\]
From which we find $Ax+A+B=x-1$ or $A=1$ and $B=-2$. 

\mynote{\textbf{method of residues.} The coefficients $A$ and $B$ can also be found
in the following way.
If I multiply equation~\eqref{eq:2888323}
by $(x+1)^2$ and then let $x$ tend to $-1$,
the coefficient $A$ on the right-hand side disappears 
and we immediately obtain $B=-2$.
Now considering the fact that the term 
with the coefficient $B$ compensates the first order of infinity 
of the term on the left-hand side (formally: bring to the left 
of the equation the term with coefficient $B$)
if we multiply the equation by $x+1$ and then let $x$ tend to $-1$
Now subtracting $-2/(x+1)^2$ from both sides eliminates 
the coefficient $B$ and we immediately obtain $A=1$.}

Thus 
\[
  \int \frac{x-1}{(x+1)^2}\, dx
  = \int \frac{1}{x+1}\, dx - 2 \int \frac{1}{(x+1)^2}
  \ni \ln \abs{x+1} + \frac{2}{x+1}. 
\]
\end{proof}

\begin{example}[denominator of second degree irreducible]
Calculate 
\[
 \int \frac{x+2}{x^2+2x+3}\, dx
\]
\end{example}
\begin{proof}[Solution.]
In this case, the denominator has no real roots (the discriminant is negative)
and therefore cannot be factored.
One must then proceed in two steps: 
\begin{enumerate}
  \item eliminate the term with $x$ in the numerator by subtracting 
  an appropriate derivative of the denominator (thus obtaining the derivative 
  of the logarithm of the denominator);
  \item complete the square in the denominator to reduce to the 
  derivative of the arctangent. 
\end{enumerate} 
For the first step, observe that $(x^2+2x+3)' = 2x+2$ and dividing the numerator
by $2x+2$ gives $x+2 = \frac 1 2 (2x+2) + 1$ from which:
\[
\int \frac{x+2}{x^2+2x+3} \, dx
= \frac 1 2 \int \frac{2x + 2}{x^2+2x+3}\, dx + \int \frac{1}{x^2+2x+3}\, dx. 
\]
But the first integral is immediate 
\[
\int \frac{2x+2}{x^2+2x+3}\, dx \ni \ln (x^2+2x+3).
\]
For the second integral, the idea is to complete the 
square to eliminate the linear term and then 
reduce with a change of variable 
$y=(x+1)/\sqrt 2$, $dx = \sqrt 2 dy$ 
to $\int \frac{1}{1+y^2}\, dy\ni \arctg y$:
\[
  \int \frac{1}{x^2+2x+3}\, dx 
  = \int \frac{1}{\enclose{x+1}^2 + 2}\, dx 
  = \frac 1 2 \int \frac{1}{\enclose{\frac{x+1}{\sqrt 2}}^2+1}\, dx 
  \ni \frac{\sqrt 2}{2} \arctg \frac{x+1}{\sqrt 2}.
\]
In conclusion, the sought integral is:
\[
  \frac 1 2 \ln (x^2+2x+3) + \frac{\sqrt 2}{2} \arctg \frac{x+1}{\sqrt 2}.
\]
\end{proof}

\subsection{Example of the General Method}

The general method can be intuited with the following 
example (decidedly more complicated).

\begin{example}
  Calculate
  \[
   \int \frac{x^{10}+x^8-x^6+1}{x^8-x^7+2x^6-2x^5+x^4-x^3}\, dx
  \]
\end{example}
\begin{proof}[Solution.]
First, we want to reduce to a rational function
where the degree of the polynomial in the numerator is less
than the degree of the polynomial in the denominator. To do this, we perform
the division between polynomials
(theorem~\ref{th:divisione_polinomi}):
\[
  \frac{x^{10}+x^8-x^6+1}{x^8-x^7+2x^6-2x^5+x^4-x^3}
  = x^2 + x + \frac{x^4+1}{x^8-x^7+2x^6-2x^5+x^4-x^3}.
\]
The integral of $x^2+x$ is immediate, so we can focus
on the rational function with numerator $x^4+1$.
To proceed, it is now necessary to determine all the roots, real and complex, of
the polynomial in the denominator so that it can be factored
as a product of linear and quadratic factors
(theorem~\ref{th:fattorizzazione_polinomio_reale}).
There is no general method for finding the roots of a polynomial.
In this particular case, it is observed
that the polynomial is divisible by $x^3$. Then it is observed that the
polynomial
of degree five resulting is zero for $x=1$ (fortunate case!) and
is therefore divisible by $x-1$ (Ruffini's theorem~\ref{th:Ruffini}).
The polynomial of degree four
is finally a perfect square, and we obtain the sought factorization:
\[
x^8-x^7+2x^6-2x^5+x^4-x^3 = x^3(x-1)(x^2+1)^2.
\]
Now Hermite's decomposition theorem (theorem~\ref{th:Hermite} which we will
later prove)
guarantees that every rational function of the form
\[
\frac{P(x)}{x^3(x-1)(x^2+1)^2}
\]
with $\deg P < 8$ (the degree of the denominator) can be
written in the following way, with
appropriate constants $A,B,C,\dots,H$:
\begin{equation}\label{eq:438943}
  \frac{A}{x} + \frac{B}{x-1} + \frac{Cx+D}{x^2+1}
  + \enclose{\frac{Ex^3 + Fx^2 + G x + H}{x^2(x^2+1)}}'.
\end{equation}
The first terms have in the denominator the factors into which we have
decomposed the denominator of the original function. In the numerator
there is always a polynomial of degree less than the denominator
and thus there is a constant above the linear factors and a linear function
above the quadratic terms. Then we add the derivative of a function
rational where the denominator has the same factors but raised to a
power one less than what was in the original polynomial.
In the numerator, again, a generic polynomial of degree less than the
denominator.
To find the coefficients $A,B,C, \dots, H$ 
(the number of coefficients is equal to the degree of the polynomial in the
denominator)
it is sufficient to perform the derivative in~\eqref{eq:438943}, 
put all terms over a common denominator 
and equate with the original rational function:
\mynote{%
When performing the derivative of the rational function, it is convenient
to consider the rational function as the product of the numerator
by the reciprocals of the factors in the denominator and use the formula
for the derivative of a product (instead of the formula for the derivative of
the quotient)%
}
\begin{align*}
  &\frac{A}{x} + \frac{B}{x-1} + \frac{Cx+D}{x^2+1}
  + \frac{3 E x^2 + 2 F x + G}{x^2(x^2+1)} \\
    & \quad + (E x^3 + F x^2 + G x + H)\cdot \left[\frac{-2}{x^3}\cdot
  \frac{1}{x^2+1}
  + \frac{1}{x^2}\cdot\frac{-2x}{(1+x^2)^2}\right]
\end{align*}
multiplying everything by $x^3(x-1)(x^2+1)^2$ we obtain
\begin{align*}
  &Ax^2(x-1)(x^2+1)^2 + Bx^3(x^2+1)^2 \\
  &\quad + (Cx+D)x^3(x-1)(x^2+1) + (3Ex^2+2Fx+G)x(x-1)(x^2+1)\\
  &\quad + (Ex^3+Fx^2+Gx+H)(-2(x-1)(x^2+1)-2x\cdot x(x-1))
\end{align*}
or
\begin{align*}
  &(A+B+C)x^7+(-A-C+D-E)x^6 + (2A+2B+C-D+E-2F)x^5 \\
  &\quad  +(-2A-C+D+E+2F-3G)x^4 + (A+B-D-E+3G-4H)x^3 \\
  &\quad  + (-A-G+4H)x^2 + (G-2H) x + 2H.
\end{align*}
Now we impose that this last polynomial is identically equal
to the original polynomial $x^4+1$ obtaining the following linear system
\[
  \begin{dcases}
    \scriptstyle
    A+B+C=0 \\[-1ex]
    \scriptstyle
    -A-C+D-E=0 \\[-1ex]
    \scriptstyle
    2A + 2B + C - D + E - 2 F = 0\\[-1ex]
    \scriptstyle
    -2A-C+D+E+2F-3G = 1\\[-1ex]
    \scriptstyle
    A+B-D-E+3G-4H=0\\[-1ex]
    \scriptstyle
    -A-G+4H=0\\[-1ex]
    \scriptstyle
    G-2H=0\\[-1ex]
    \scriptstyle
    2H=1
  \end{dcases}
\]
which, when solved, gives us: $A=1$, $B=\frac 1 2$, $C = -\frac 3 2$, $D=1$,
$E=\frac 3 2$, $F=1$, $G=1$, $H=\frac 1 2$.
We have obtained, ultimately:
\begin{align*}
\frac{x^4+1}{x^3(x-1)(x^2+1)^2}
&= \frac{1}{x} + \frac{\frac 1 2}{x-1} + \frac{1 - \frac 3 2 x}{x^2+1}
+ \enclose{\frac{\frac 3 2 x^3 + x^2 + x + \frac 1 2}{x^2(x^2+1)}}'
\end{align*}
and observing that
\[
 \int \frac{1 - \frac 3 2 x}{x^2+1}\, dx
 = \int \frac{1}{x^2+1}\, dx - \frac 3 4 \int \frac{2x}{x^2+1}\, dx
 \ni \arctg x - \frac 3 4 \ln(x^2+1)
\]
we finally obtain
\begin{align*}
\int \frac{x^4+1}{x^3(x-1)(x^2+1)^2}\, dx
 &\ni \ln \abs{x} + \frac 1 2 \ln\abs{x-1} + \arctg x - \frac 3 4 \ln(x^2+1) \\
 &\quad + \frac{\frac 3 2 x^3 + x^2+x+\frac 1 2}{x^2(x^2+1)}
\end{align*}
\end{proof}

The method outlined in the previous example is an algorithm 
valid in general:
\index{algorithm!for the integration of rational functions}%
\index{integration!of rational functions!algorithm}%
\index{function!rational!integration}%
\index{decomposition!into partial fractions}%
\index{decomposition!into partial fractions}%
\begin{enumerate}
  \item Perform the division with remainder to reduce
  to the case where the numerator of the rational function
  has a degree lower than the denominator.
  \item Decompose the denominator as a product of
  linear and quadratic polynomials.
  \item Use Hermite's decomposition (theorem~\ref{th:Hermite})
  to write the rational function as a sum of \emph{partial fractions}
  plus the derivative of another rational function. The coefficients
  of the decomposition are found by solving a linear system.
  \item Integrate the partial fractions.
\end{enumerate}

To conclude the previous algorithm, it is necessary to know how to integrate the
partial fractions. If the denominator is linear, the integral is immediate
(resulting in a logarithm).

If the denominator is quadratic, the integral can be slightly more complicated,
and it is useful to know how to proceed. The quadratic terms
will be of the form:
\begin{equation}\label{eq:0943784}
\frac{Ax + B}{x^2 + \alpha x + \beta}
\end{equation}
with $\alpha^2-4\beta<0$ (otherwise, the quadratic polynomial
in the denominator can be decomposed into two linear factors, using
the quadratic formula~\eqref{eq:secondo_grado}). The idea is to perform a linear
change
of variable to reduce to a denominator
of the form $y^2+1$. To do this, use the method of completing the square as we
have already done in~\eqref{eq:24589}:
\[
  x^2 + \alpha x + \beta
  = \enclose{x+\frac\alpha 2}^2  + \beta - \frac{\alpha^2}{4}
  = \dots
\]
at this point, factor out $\beta-\frac{\alpha^2}{4}$
(which by hypothesis is positive)
and bring it inside the square:
\[
 \dots = \enclose{\beta-\frac{\alpha^2}{4}}\cdot\Enclose{\enclose{\frac{x+\frac \alpha 2}{\sqrt{\beta -\frac{\alpha^2}{4}}}}^2 + 1}
\]
and thus through the substitution $y=\frac{x + \frac \alpha 2 }{\sqrt{\beta -
\frac{\alpha^2}{4}}}$
we have reduced our rational function~\eqref{eq:0943784}
to the form
\[
  \frac{a y + b}{y^2 + 1} = \frac a 2 \frac{2y}{y^2+1} + b\frac{1}{y^2+1}
\]
which can be integrated immediately.

\subsection{Partial Fraction Decomposition and Hermite Decomposition}
\index{decomposition!into partial fractions}%
\index{decomposition!into partial fractions}%
\index{decomposition!of Hermite}%
\index{decomposition!of Hermite}%

In the following, we propose the statements and proofs that guarantee
the functioning of the algorithm for the integration of rational functions.
These proofs can hardly be completed
without using a minimum of properties of complex functions. In particular, it is
necessary to know that the derivatives of elementary complex functions
(sums, products, quotients, composition) are carried out exactly
with the same rules that apply to real functions (see
chapter~\ref{sec:derivata_complessa}).

\begin{lemma}[independence of partial fractions]
\label{lemma:72995}
Let $\lambda_1,$ $\dots,$ $\lambda_n \in \CC$ be distinct complex numbers
and let $p_1,\dots,p_n \in \NN\setminus\ENCLOSE{0}$.
Let $\Omega = \CC \setminus \ENCLOSE{\lambda_1, \dots, \lambda_n}$.
Then the functions
$u_{\lambda_k}^{j}\colon \Omega\subset \CC \to \CC$
\begin{equation}
\label{eq:483818}
  u_{\lambda_k}^{j}(z) = \frac{1}{(z-\lambda_k)^{j}}
  \qquad k=1,\dots,n, \quad j=1,\dots, p_k
\end{equation}
are linearly independent elements of the complex vector space
$\CC^\Omega$.
\end{lemma}
%
\begin{proof}
We proceed by induction on the number of functions
$N=p_1+ \dots + p_n$.
If $N=1$, we have a single function that is not identically
zero (in fact: it never vanishes)
and thus constitutes a linearly independent set.

Suppose then we have a set consisting of more functions
in the form~\eqref{eq:483818}.
and have
a complex linear combination identically zero:
\[
\sum_{k=1}^n \sum_{j=1}^{p_k} \alpha_{k,j} u_{\lambda_k}^j = 0.
\]
Consider the function $f\colon \Omega \to \CC$
defined by
\begin{equation}
\label{eq:567384}
  f(z)
    = (z-\lambda_n)^{p_n}\sum_{k=1}^n \sum_{j=1}^{p_k}
  \frac{\alpha_{k,j}}{(z-\lambda_k)^{j}}
    = \sum_{k=1}^n \sum_{j=1}^{p_k}
  \alpha_{k,j}\frac{(z-\lambda_n)^{p_n}}{(z-\lambda_k)^{j}}.
\end{equation}
This function is identically zero
in $\Omega$ as it is obtained by multiplying the identically
zero linear combination by the factor $(z-\lambda_n)^{p_n}$.
But it is observed that we have
\[
  \lim_{z\to \lambda_n} f(z) = \alpha_{n,p_n}
\]
since if $k\neq n$ we have
\[
\lim_{z\to \lambda_n}\frac{(z-\lambda_n)^{p_n}}{(z-\lambda_k)^{p_k}} = 0
\]
since $\lambda_k\neq \lambda_n$ and the numerator
tends to zero, if instead $k=n$ and $j<p_n$ we still have
\[
\lim_{z\to \lambda_n}\frac{(z-\lambda_n)^{p_n}}{(z-\lambda_k)^{j}}
= \lim_{z\to \lambda_n}(z-\lambda_n)^{p_n-j} = 0.
\]
Thus $\alpha_{n,p_n}=0$.
But then we have a zero linear combination of $N-1$
functions and by the inductive hypothesis, we can conclude
that all other coefficients are also zero.
\end{proof}

\begin{theorem}[complex partial fraction decomposition]
\label{th:fratti_semplici_complessi}%
\index{decomposition!into complex partial fractions}%
\index{decomposition!into complex partial fractions}%
\index{partial fractions!complex decomposition}%
If $\lambda_1, \dots, \lambda_n$ are the distinct complex roots
of the polynomial $Q$ with complex coefficients,
and $p_1,\dots, p_n$ are the respective multiplicities
with $p_1 + \dots + p_n = \deg Q$
and if $P$ is any polynomial with complex coefficients
with $\deg P < \deg Q$,
then there exist unique polynomials $R_1, \dots, R_n$
with complex coefficients with $\deg R_k < p_k$ such that
\begin{equation}\label{eq:46772341}
  \frac{P(z)}{Q(z)}
  = \sum_{k=1}^n \frac{R_k(z)}{(z-\lambda_k)^{p_k}}
  \qquad \forall z \in \CC\setminus\ENCLOSE{\lambda_1,\dots,\lambda_n}.
\end{equation}
\end{theorem}
%
\begin{proof}
Let $N=\deg Q$ and consider
the set of polynomials with complex coefficients
of degree less than $N$:
\[
  V_N = \ENCLOSE{P\in \CC[z]\colon \deg P < N}.
\]
This is a vector subspace of
$V=\CC^\Omega$ and has dimension
$N$ (a basis is given by the monomials $1, z, z^2, \dots, z^{N-1}$).
Consequently, it is easy to verify that also the vector space
\[
  W_Q = V_N \cdot \frac{1}{Q} = \ENCLOSE{\frac{P}{Q}\colon \text{$\deg P < N$}}
\]
whose elements are the rational functions on the left side
of~\eqref{eq:46772341},
has the same dimension $\dim W_Q = \dim V_N = N$.
On the other hand, the right side of~\eqref{eq:46772341}
is a linear combination of functions
$u_{\lambda_k}^j = 1/(z-\lambda_k)^j$ which are
also elements of $W_Q$ since $(z-\lambda_k)^j$ divides $Q(x)$
if $j\le p_k$.
But by lemma~\ref{lemma:72995} the vectors $u_{\lambda_k}^j$ are
$N$ independent vectors and thus
are a basis of $W$.
This means that every element of $W$ can
be uniquely written as a linear combination of the
$u_{\lambda_k}^j$, that is, for every
polynomial $P$ with $\deg P < \deg Q$ we have
\begin{equation}\label{eq:458934}
 \frac{P(z)}{Q(z)}
 = \sum_{k=1}^n \sum_{j=1}^{p_k} \frac{A_{k,j}}{(z-\lambda_k)^j}
  = \sum_{k=1}^n \frac{\displaystyle \sum_{j=1}^{p_n} A_{k,j}\cdot
 (z-\lambda_k)^{p_n-j}}{(z-\lambda_k)^{p_n}}
\end{equation}
and setting
\[
  R_k(z) = \sum_{j=1}^{p_n} A_{k,j}\cdot(z-\lambda_k)^{p_n-j}
\]
it is clear that $R_k$ is a polynomial of degree less than $p_n$
uniquely determined by the coefficients $A_{k,j}$.
\end{proof}


\begin{theorem}[Hermite decomposition]
\label{th:Hermite}%
\index{Hermite!decomposition into partial fractions}%
\index{theorem!of Hermite}%
Let $Q(x)$ be a monic polynomial with real coefficients
with real roots $\lambda_1, \dots, \lambda_n$ of multiplicity
$p_1, \dots, p_n$ and non-real complex conjugate roots
$\mu_1, \bar \mu_1, \dots, \mu_m, \bar \mu_m$ of multiplicity
$q_1, \dots, q_m$. Setting $\alpha_k = 2\Re \mu_k$ and $\beta_k = \abs{\mu_k}^2$
we can write (thanks to theorem~\ref{th:fattorizzazione_polinomio_reale})
\begin{equation}\label{eq:8845638}
    Q(x) = \prod_{k=1}^n (x-\lambda_k)^{p_k} \cdot \prod_{k=1}^m (x^2+ \alpha_k
  x + \beta_k)^{q_k}.
\end{equation}
Let $P(x)$ be any polynomial with real coefficients
with $\deg P < \deg Q$.

Then setting
\[
 \tilde Q(x) = \prod_{k=1}^n (x-\lambda_k)^{p_k-1} \cdot \prod_{k=1}^m (x^2+\alpha_k x+\beta_k)^{q_k-1}
\]
there exist (unique) real coefficients
$A_1,\dots, A_n$, $B_1, \dots, B_m$, $C_1,\dots, C_m$
and a polynomial with real coefficients $R(x)$ with $\deg R < \deg \tilde Q$
such that:
\begin{equation}\label{eq:scomposizione_hermite}
  \frac{P(x)}{Q(x)} = \sum_{k=1}^n \frac{A_k}{x-\lambda_k}
  + \sum_{k=1}^m \frac{B_k x + C_k}{x^2+\alpha_k x + \beta_k}
  + \enclose{\frac{R(x)}{\tilde Q(x)}}'.
\end{equation}
By performing the derivative, the right side can be
rewritten with a common denominator $Q(x)$
and by equating the numerators of the right and left sides
a linear system with $\deg Q$ unknowns will be obtained
which will have a unique solution.
\end{theorem}
%
\begin{proof}
Consider the functions
\begin{align*}
   u_\lambda^j(z) &= \frac{1}{(z-\lambda)^j}\\
   v_\mu(z) &= \frac{1}{z^2-2 \Re \mu\cdot z + \abs{\mu}^2},\\
   w_\mu(z) &= \frac{z}{z^2-2 \Re \mu\cdot z + \abs{\mu}^2},\\
   t_j(z) &= \frac{z^j}{\tilde Q(z)}.
\end{align*}
Let $\lambda_1, \dots, \lambda_N$ be all the complex roots
of the polynomial $Q(z)$ with $p_1, \dots, p_N$
the respective multiplicities.
Thanks to lemma~\ref{lemma:72995}, we can assert,
as we did in the proof of
theorem~\ref{th:fratti_semplici_complessi}, that the
vector space $V_Q$ of all rational functions of the form $\frac{P(z)}{Q(z)}$
with $Q$ fixed and $\deg P<\deg Q$
is generated by the linear combinations
with complex coefficients of the functions
\[
  u_{\lambda_k}^j, \qquad k=1, \dots, N, \quad j=1, \dots, p_k.
\]
Similarly, the vector space $V_{\tilde Q}$
is generated by the functions
\[
  u_{\lambda_k}^j, \qquad k=1, \dots, N, \quad j=1, \dots, p_k-1.
\]
Now observe that for $j>1$ we have
\[
 \frac {d}{dz} u_\lambda^{j-1} =
 \frac {d}{dz} \frac{1}{(z-\lambda)^{j-1}}
 = \frac{1-j}{(z-\lambda)^j} = (1-j) u_\lambda^j
\]
which means that the linear operator $\frac d {dz}$
maps $V_{\tilde Q}$ into $V_Q$, more precisely
the image of $\frac{d}{dz} V_{\tilde Q}$ is the subspace
of $V_Q$ generated by the functions $u_{\lambda_k}^j$
with $j>1$. Such an operator
is also injective as it maps vectors of
the basis of $V_{\tilde Q}$ into vectors of the basis of $V_Q$
which are certainly independent.
We thus have a decomposition into a direct sum:
\[
  V_Q = V_1 \oplus \frac{d}{dz} V_{\tilde Q}
\]
where $V_1$ is the space generated
by the functions $u_{\lambda_1}^1, \dots, u_{\lambda_N}^1$.

But also the functions $t_0, \dots, t_{M-1}$ with $M=\deg \tilde Q$
are a basis of $V_{\tilde Q}$, thus a basis
of $V_Q$ is given by
\[
  u_{\lambda_1}, \dots, u_{\lambda_N},
  t_0', \dots, t_{M-1}'
\]
where $t_m'$ is the derivative of $t_m$.
Therefore, since every linear combination
of $u_{\mu}$ and $u_{\bar \mu}$ can be written
as a linear combination of $v_\mu$ and $w_\mu$
we obtain that another basis of $V_Q$ is
\begin{equation}\label{eq:63928}
  u_{\lambda_1}, \dots, u_{\lambda_n},
  v_{\mu_1}, \dots, v_{\mu_m},
  w_{\mu_1}, \dots, w_{\mu_m},
  t_0', \dots, t_{M-1}'.
\end{equation}
For brevity, we call $e_1, \dots, e_N$
the basis of $V_Q$ listed in~\eqref{eq:63928}.
Each of these functions can be written as
a ratio of polynomials with real coefficients
whose denominator is a divisor of the polynomial $Q$
which is also a polynomial with real coefficients.
Thus the functions $Q\cdot e_1, \dots, Q\cdot e_N$
are actually polynomials with real coefficients
and must necessarily be a basis of the complex vector space $Q\cdot V_Q$, which
is nothing but the set
of all polynomials with complex coefficients
of degree less than $N$. In particular, these polynomials
are independent.
But then, viewed as polynomials with
real coefficients, they are also independent and are thus
a basis of the space of all polynomials with real coefficients
of degree less than $N$.
Therefore, there must exist real coefficients $\alpha_1, \dots, \alpha_N$
such that we have
\[
  P(x) = \sum_{k=1}^N \alpha_k\cdot Q(x)\cdot e_k(x)
\]
from which
\[
  \frac{P(x)}{Q(x)} = \sum_{k=1}^N \alpha_k \cdot e_k(x).
\]
Recalling that the functions $e_k$ are nothing but the functions
listed in~\eqref{eq:63928}, we finally obtain the sought decomposition.
\end{proof}